\section{September 4, 2026}

We continue the study of Runge--Kutta methods, first constructing
families of explicit methods and then introducing embedded methods and
stability functions. We use the notation established in the previous
lectures: \(\tau_n=t_{n+1}-t_n\), \(x_n\) is the numerical solution,
and \(\Phi\) and \(\Psi\) denote the exact and numerical flows.

\subsection{Runge--Kutta methods and order conditions}

Recall that an \(s\)-stage Runge--Kutta method is
\begin{align*}
  k_i
  &=f\left(t_n+c_i\tau_n,
    x_n+\tau_n\sum_{j=1}^s a_{ij}k_j\right),
    \qquad i=1,\ldots,s,\\
  x_{n+1}
  &=x_n+\tau_n\sum_{i=1}^s b_i k_i.
\end{align*}
Here \(s\) is the number of stages, \(A=(a_{ij})\) contains the stage
coefficients, \(c_i\) is the \(i\)-th stage node, and \(b_i\) is the
\(i\)-th quadrature weight. The internal consistency relation is
\[
  c=A\mathbf{1}.
\]
It makes the method agree on the original nonautonomous equation and
its autonomous augmentation. The method is explicit precisely when
\(A\) is strictly lower triangular.

For the autonomous equation \(x'=f(x)\), a method has order \(p\) when
\[
  \bigl\lVert\Phi^\tau(x)-\Psi^\tau(x)\bigr\rVert
  =O(\tau^{p+1})
  \qquad\text{as }\tau\to0.
\]
The first four order conditions, derived in the previous lecture, are
\begin{align*}
  b^{\mathsf T}\mathbf{1}&=1,
  &b^{\mathsf T}c&=\frac12,\\
  b^{\mathsf T}c^{\circ2}&=\frac13,
  &b^{\mathsf T}Ac&=\frac16.
\end{align*}
They are sufficient through order three. To obtain order four, four
more conditions are needed:
\begin{align*}
  b^{\mathsf T}c^{\circ3}&=\frac14,
  &b^{\mathsf T}\bigl(c\circ(Ac)\bigr)&=\frac18,\\
  b^{\mathsf T}A(c^{\circ2})&=\frac1{12},
  &b^{\mathsf T}A^2c&=\frac1{24},
\end{align*}
where \(\circ\) denotes componentwise multiplication. Thus one, two,
four, and eight conditions must be satisfied for orders one, two,
three, and four, respectively.

Forward Euler uses one stage and has order one, while the explicit
trapezoidal method uses two stages and has order two. We next continue
the construction of higher-order explicit methods.

\subsection{Explicit three-stage methods of order three}

For an explicit three-stage method, internal consistency gives the
tableau
\[
  \begin{array}{c|ccc}
    0&0&0&0\\
    c_2&a_{21}&0&0\\
    c_3&a_{31}&a_{32}&0\\ \hline
      &b_1&b_2&b_3
  \end{array},
  \qquad
  c_2=a_{21},
  \quad
  c_3=a_{31}+a_{32}.
\]
There are six free coefficients and four order conditions, so the
third-order methods generically form a two-parameter family. Let
\[
  c_2=\alpha,
  \qquad
  c_3=\beta.
\]
Solving the four order conditions gives
\[
  \begin{array}{c|ccc}
    0&0&0&0\\[2pt]
    \alpha&\alpha&0&0\\[4pt]
    \beta
      &\dfrac{\beta\bigl(3\alpha(1-\alpha)-\beta\bigr)}
              {\alpha(2-3\alpha)}
      &\dfrac{\beta(\beta-\alpha)}{\alpha(2-3\alpha)}
      &0\\[8pt] \hline
      &\dfrac{6\alpha\beta-3\alpha-3\beta+2}{6\alpha\beta}
      &\dfrac{3\beta-2}{6\alpha(\beta-\alpha)}
      &\dfrac{2-3\alpha}{6\beta(\beta-\alpha)}
  \end{array},
\]
provided the displayed denominators are nonzero.

A useful subfamily has \(\beta=1\). Its weights are
\[
  b_1=\frac{3\alpha-1}{6\alpha},
  \qquad
  b_2=\frac{1}{6\alpha(1-\alpha)},
  \qquad
  b_3=\frac{2-3\alpha}{6(1-\alpha)}.
\]
All three weights are nonnegative when
\(\alpha\in[1/3,2/3)\), but positivity of the weights alone does not
guarantee that one method is better than another.

\begin{example}[Kutta's third-order method]
  Taking \(\alpha=1/2\) and \(\beta=1\) gives
  \[
    \begin{array}{c|ccc}
      0&0&0&0\\
      \tfrac12&\tfrac12&0&0\\
      1&-1&2&0\\ \hline
       &\tfrac16&\tfrac23&\tfrac16
    \end{array}.
  \]
  Thus
  \begin{align*}
    k_1&=f(t_n,x_n),\\
    k_2&=f\left(t_n+\frac{\tau_n}{2},
      x_n+\frac{\tau_n}{2}k_1\right),\\
    k_3&=f\left(t_n+\tau_n,
      x_n+\tau_n(-k_1+2k_2)\right),\\
    x_{n+1}&=x_n+\frac{\tau_n}{6}(k_1+4k_2+k_3).
  \end{align*}
\end{example}

\subsection{Explicit four-stage methods of order four}

Before imposing internal consistency, a general explicit four-stage
tableau contains thirteen coefficients: three nontrivial nodes, six
strictly lower-triangular entries of \(A\), and four weights. Internal
consistency removes three degrees of freedom, and the eight
fourth-order conditions generically leave two free parameters.

In the usual nonconfluent family one may take
\[
  c_1=0,
  \qquad
  c_2=\alpha,
  \qquad
  c_3=\beta,
  \qquad
  c_4=1,
\]
with distinct stage nodes. The generic formulas exclude degenerate
parameter choices; one convenient nondegeneracy condition is
\[
  \alpha\beta(\alpha-1)(\beta-1)(\alpha-\beta)(2\alpha-1)
  \bigl(6\alpha\beta-4\alpha-4\beta+3\bigr)\neq0.
\]

\begin{example}[The three-eighths rule]
  Choosing \(\alpha=1/3\) and \(\beta=2/3\) produces the
  fourth-order three-eighths rule
  \[
    \begin{array}{c|cccc}
      0&0&0&0&0\\
      \tfrac13&\tfrac13&0&0&0\\
      \tfrac23&-\tfrac13&1&0&0\\
      1&1&-1&1&0\\ \hline
       &\tfrac18&\tfrac38&\tfrac38&\tfrac18
    \end{array}.
  \]
  Its stages are
  \begin{align*}
    k_1&=f(t_n,x_n),\\
    k_2&=f\left(t_n+\frac{\tau_n}{3},
      x_n+\frac{\tau_n}{3}k_1\right),\\
    k_3&=f\left(t_n+\frac{2\tau_n}{3},
      x_n+\tau_n\left(-\frac13k_1+k_2\right)\right),\\
    k_4&=f\left(t_n+\tau_n,
      x_n+\tau_n(k_1-k_2+k_3)\right),\\
    x_{n+1}&=x_n+\frac{\tau_n}{8}
      (k_1+3k_2+3k_3+k_4).
  \end{align*}
\end{example}

Repeated nodes give a separate, \emph{confluent} family. Setting
\(c_2=c_3=1/2\) and letting \(\delta=a_{32}\neq0\) yields
\begin{equation}\label{eq:confluent-rk4-family-sep-04}
  \begin{array}{c|cccc}
    0&0&0&0&0\\[2pt]
    \tfrac12&\tfrac12&0&0&0\\[2pt]
    \tfrac12&\tfrac12-\delta&\delta&0&0\\[4pt]
    1&0&1-\dfrac{1}{2\delta}&\dfrac{1}{2\delta}&0\\[6pt] \hline
     &\tfrac16&\dfrac{4\delta-1}{6\delta}
       &\dfrac{1}{6\delta}&\tfrac16
  \end{array}.
\end{equation}

\begin{example}[Classical RK4]
  The choice \(\delta=1/2\) in
  \eqref{eq:confluent-rk4-family-sep-04} gives the classical tableau
  \[
    \begin{array}{c|cccc}
      0&0&0&0&0\\
      \tfrac12&\tfrac12&0&0&0\\
      \tfrac12&0&\tfrac12&0&0\\
      1&0&0&1&0\\ \hline
       &\tfrac16&\tfrac13&\tfrac13&\tfrac16
    \end{array}.
  \]
  Equivalently,
  \begin{align*}
    k_1&=f(t_n,x_n),\\
    k_2&=f\left(t_n+\frac{\tau_n}{2},
      x_n+\frac{\tau_n}{2}k_1\right),\\
    k_3&=f\left(t_n+\frac{\tau_n}{2},
      x_n+\frac{\tau_n}{2}k_2\right),\\
    k_4&=f(t_n+\tau_n,x_n+\tau_n k_3),\\
    x_{n+1}&=x_n+\frac{\tau_n}{6}
      (k_1+2k_2+2k_3+k_4).
  \end{align*}
\end{example}

\subsection{Embedded pairs}

Adaptive solvers need local error estimates. They obtain them by
combining the same stages in two different ways. Given fixed \(A\) and \(c\), choose two
weight vectors \(b\) and \(\widehat b\) and form
\begin{align*}
  x_{n+1}^{[p]}
  &=x_n+\tau_n\sum_{i=1}^s b_i k_i,\\
  x_{n+1}^{[p+1]}
  &=x_n+\tau_n\sum_{i=1}^s \widehat b_i k_i.
\end{align*}
The two formulas share all evaluations \(k_i\), so
\[
  \operatorname{err}_n
  =\bigl\lVert x_{n+1}^{[p+1]}-x_{n+1}^{[p]}\bigr\rVert
\]
is an inexpensive local error estimate. If this quantity is well below
the requested tolerance, the next step size can be increased. If it is
too large, the step size is decreased and the step may be rejected.

MATLAB's \texttt{ode45} uses a seven-stage Dormand--Prince embedded
pair of orders four and five. The minimal number of stages for an
explicit Runge--Kutta method grows with its order; the low-order values
discussed in class are
\[
  \begin{array}{c|cccccccc}
    p&1&2&3&4&5&6&7&8\\ \hline
    s_{\min}&1&2&3&4&6&7&9&11
  \end{array}.
\]

\subsection{The stability function}

Apply a Runge--Kutta method to the linear test equation
\begin{equation}\label{eq:linear-test-equation-sep-04}
  y'(t)=\lambda y(t).
\end{equation}
With \(z=\tau_n\lambda\), one numerical step has the form
\begin{equation}\label{eq:stability-update-sep-04}
  y_{n+1}=R(z)y_n,
\end{equation}
where \(R\) is the \emph{stability function}. The exact amplification
factor is \(e^z\).

For a Runge--Kutta method, the stage equations applied to
\eqref{eq:linear-test-equation-sep-04} give
\[
  (I-zA)k=\lambda y_n\mathbf{1}.
\]
Substituting the resulting stages into the update gives
\begin{equation}\label{eq:rk-stability-function-sep-04}
  R(z)=1+z b^{\mathsf T}(I-zA)^{-1}\mathbf{1}.
\end{equation}
If the method has order \(p\), then necessarily
\begin{equation}\label{eq:stability-order-condition-sep-04}
  R(z)-e^z=O(z^{p+1})
  \qquad\text{as }z\to0.
\end{equation}
For example, the eight fourth-order conditions imply
\[
  R(z)=1+z+\frac{z^2}{2!}+\frac{z^3}{3!}
       +\frac{z^4}{4!}+O(z^5).
\]
Condition \eqref{eq:stability-order-condition-sep-04} is necessary for
order \(p\), but it only tests a linear ODE and is not sufficient to
establish order \(p\) for general nonlinear equations.
